\documentclass{article}
\usepackage{graphicx} % Required for inserting images
\usepackage[margin=1in]{geometry}
\usepackage{amsmath}
\usepackage{amsfonts}
\usepackage{amsthm}
\usepackage{thm-restate}
\usepackage[]{hyperref}
\usepackage{cleveref}
\usepackage{xcolor}
\usepackage{bbm}
\usepackage{dirtytalk}
\usepackage{graphicx}
\usepackage{subcaption}
\usepackage{float}
\usepackage{tikz}
\usepackage{algorithm}
\usepackage[noend]{algpseudocode}
\usetikzlibrary{arrows.meta,decorations.markings, calc, positioning}


\declaretheorem[name=Theorem,numberwithin=section]{thm}
\declaretheorem[name=Lemma,numberwithin=section]{lemma}
\declaretheorem[name=Assumption,numberwithin=section]{asmpt}
\declaretheorem[name=Proposition,numberwithin=section]{prop}
\declaretheorem[name=Corollary,numberwithin=section]{corollary}
\declaretheorem[name=Remark,numberwithin=section]{rmk}
\declaretheorem[name=Definition,numberwithin=section]{defn}

\newcommand{\man}{\mathcal{M}}

\definecolor{cyan}{cmyk}{1, 0.4, 0, 0} 
\newcommand{\ndfg}[1]{{\color{magenta}[Florian: #1]}}
\newcommand{\tls}[1]{{\color{cyan}[TLS: #1]}}


\usepackage[numbers]{natbib}
\bibliographystyle{unsrtnat}

\title{Parallel Transport on HK$(\mathbb{R}^d)$}
\author{Florian Gunsilius, Gonzalo Mena, Tristan Saidi, Larry Wasserman}

\begin{document}

\maketitle


\section{The 2-Wasserstein and HK tangent structure}
\textcolor{cyan}{(TLS: I'm going to re-write the tangent space structure in the $W_2$ case to make it concrete for myself.)} A curve in $\mathcal{P}_2(\mathbb{R}^d)$ is a family of densities indexed by time $\{\rho_t\}_{t\in [0,\tau]}$, $T_t: \mathbb{R}^d \times [0, \tau] \rightarrow \mathbb{R}$ where $T_t$ is the map $(x,t) \mapsto \rho_t$. This curve can be described by a velocity field $v_t(x)$, where $\rho_t$ evolves according to the continuity equation
\[\partial_t \rho_t  + \nabla \cdot (\rho_tv_t) = 0.\]
In the Wasserstein-Fisher-Rao case we relax the mass preservation constraint, yielding an augmented continuity equation given by
\[\partial_t \rho_t  + \nabla \cdot (\rho_tv_t) = 4g_t\rho_t\]
where $g_t$ describes the birth/death rate at $x \in \mathbb{R}^d$ at time $t$. Focusing on the Wasserstein case, we can characterize tangent vectors by
\[s_t \in T_{\rho_t} \mathcal{P}_2 \iff s_t = -\nabla \cdot (\rho_t v_t)\]
for some vector field $v_t: \mathbb{R}^d \rightarrow \mathbb{R}^d$. This isn't quite correct ($v_t$ needs to be the gradient of a potential function), but we'll come back to that soon. First, notice that $s_t: \mathbb{R}^d \rightarrow \mathbb{R}$ is a scalar valued function over $\mathbb{R}^d$ representing the negative divergence of the scaled vector field $\rho_tv_t$. To address why the vector field $v \equiv \nabla\varphi$: consider a vector field $w_t$ such that 
\[-\nabla \cdot (\rho_t w_t) = 0.\]
By the continuity equation, this vector field $w_t$ does not induce any change in the density anywhere (even though it moves mass). Thus, for any other tangent vector at $\rho_t$ we have 
\begin{align}
    \frac{\partial\rho_t}{\partial t} &= -\nabla \cdot (\rho_tv_t)\\  
    &= -\nabla \cdot (\rho_tv_t) + \underbrace{-\nabla \cdot(\rho_tw_t)}_{0} \\
    &= -\nabla \cdot (\rho_t(v_t +w_t)).
\end{align}
This means that $v_t$ and $v_t + w_t$ generate the same tangent vector. To establish uniqueness, we can define the equivalence relation $v_t \sim w_t$ if $-\nabla \cdot (\rho_t(v_t + w_t)) = 0$. But how do we now obtain a representative element for each equivalence class? \citet{otto2001geometry} proposes that we choose the one with minimal norm. In particular, for a tangent vector $s$ we can identify it with the vector field that solves
\begin{equation}
    v_t = \underset{s_t +\nabla\cdot (\rho_tv) = 0}{\arg\min}\int_{\mathbb{R}^d}\|v\|^2\,d\rho.
\end{equation}
By a Helmholtz decomposition $v = \nabla \varphi + w$ it becomes clear that
\begin{align}
    v_t &= \underset{s_t +\nabla\cdot (\rho_tv) = 0}{\arg\min}\int_{\mathbb{R}^d}\|v\|^2\,d\rho \\
    &= \underset{\substack{s_t +\nabla\cdot (\rho_t\nabla\varphi) = 0, \\\nabla\cdot(\rho_tw) = 0}}{\arg\min}\int_{\mathbb{R}^d}\|\nabla\varphi + w\|^2\,d\rho \\
    &= \underset{\substack{s_t +\nabla\cdot (\rho_t\nabla\varphi) = 0, \\\nabla\cdot(\rho_tw) = 0}}{\arg\min}\int_{\mathbb{R}^d}\Big[\|\nabla\varphi\|^2 + \underbrace{\langle\nabla\varphi, w\rangle}_{0} + \|w\|^2\Big]\,d\rho
\end{align}
and therefore the minimum norm solution is $v_t = \nabla\varphi$ for some potential function $\varphi: \mathbb{R}^d\rightarrow\mathbb{R}$. Now one can assign a Riemannian metric on $\mathcal{P}_2(\mathbb{R}^d)$ as follows. For vectors $u_1, u_2 \in T_{\rho}\mathcal{P}_2$ we have
\begin{equation}
    \langle u_1, u_2\rangle_{\rho} = \int_{\mathbb{R}^d}\langle\nabla\varphi_1, \nabla\varphi_2\rangle\, d\rho
\end{equation}
where $\varphi_1$ and $\varphi_2$ are the unique potential functions corresponding to $u_1$ and $u_2$ respectively. In the Wasserstein-Fisher-Rao case, we have non-uniqueness from two sources: in particular, for a given $s_t$ there are many $v_t, g_t$ satisfying the modified continuity equation
\[s_t + \nabla \cdot (\rho_tv_t) = 4g_t\rho_t.\]
One option, as suggested in \cite{vialard2015wasserstein, chizat2018interpolating, lambert2022variational}, is to choose the minimum norm solution under the Fisher-Rao metric,
\[(v_t, g_t) = \underset{\substack{s_t +\nabla\cdot (\rho_tv_t) = g_t\rho_t}}{\arg\min}\int_{\mathbb{R}^d} \left(4|g_t|^2 + \|v_t\|^2\right)\,d\rho_t.\]
As indicated by \citet{clancy2021interpolating} Proposition 27, this optimization has a unique solution 
\begin{equation}
\bar{v}_t, \bar{g}_t \in \overline{\{(\nabla\varphi, \varphi) \, | \, \varphi \in C^{\infty}(\mathbb{R}^d) \}} \label{eq: HK tangent vectors}
\end{equation}
where the closure is taken in the $L_2(\rho_t)$ sense. It follows that the HK metric tensor is
\begin{equation}
    \langle s_1, s_2 \rangle_{\rho} = \int_{\mathbb{R}^d} \left(\langle v_1, v_2\rangle + 4g_1g_2\right)\,d\rho.
\end{equation}

\section{Defining parallel transport}
\textcolor{cyan}{(TLS: again, I'm going to do this in the Wasserstein case first, and then flesh out the analogy for HK.)}
\subsection{Connections/Covariant Derivatives}
On an abstract manifold $(\man, g)$ that isn't embedded in an ambient space, we have no obvious way to compare vectors in the tangent space at two points $p, q \in \man$, $p \neq q$. Therefore, we need a way of connecting the separate vector spaces $T_p\man$ and $T_q\man$ when $p \neq q$. One such way of doing so is by defining a differential operator, which then induces a connection between tangent spaces. Note that in this section I will use the word connection and the phrase covariant derivative interchangeably -- they are fundamentally the same object.

Suppose I have vector fields $X$ and $Y$ over my manifold $\man$, and I want to consider the derivative of $Y$ in the direction of the vector field $X$ (denoted $\nabla_XY$). In the Euclidean case, this is trivial -- it amounts to the classical directional derivative,
\[\frac{d}{dt}Y(\gamma(t)) = \lim_{\Delta t\rightarrow 0}\frac{Y(\gamma(t + \Delta t)) - Y(\gamma(t))}{\Delta t}.\]
But in the general case, this is challenging precisely because we cannot subtract $Y(\gamma(t ))$ from $Y(\gamma(t + \Delta t))$. But observe the following: if I'm \textit{given} a rule $\nabla$ for differentiating vector fields, we can define parallel transport and therefore a meaningful notion of comparing tangent vectors at nearby points through the relation above (loosely speaking). This motivates the idea of the covariant derivative.
If we can define a way of differentiating vector fields against each other on $\man$ in a way that preserves the structure of the metric $g$, then we have defined a connection on $\man$.

\begin{restatable}[The Covariant Derivative]{defn}{}
    For a manifold $(\man, g)$ we define the covariant derivative $\nabla$ to be a rule that assigns to each pair of vector fields $X$, $Y$ over $\man$ another vector field $\nabla_XY$ satisfying 
    \begin{enumerate}
        \item \textbf{Linearity:} $\nabla_{aX+bZ}Y = a\nabla_XY + b\nabla_ZY.$
        \item \textbf{Leibniz/derivation property:} for $f \in C^{\infty}(\man)$ we have $\nabla_X(fY) = X(f)Y + f\nabla_XY$, where $X(f)$ is the derivative of $f$ in the direction $X$. \textcolor{red}{gm: how is this derivative defined?} \tls{see Loring Tu's An Intro to Manifolds, Proposition 8.17. One can get the directional derivative from the metric tensor with $g(X, \nabla f)$ -- the gradient of $f$, $\nabla f$ does not depend on the covariant derivative because it comes directly from the manifolds coordinate charts. }
    \end{enumerate}
\end{restatable}

While this gives us a way to connect tangent spaces, the covariant derivative might distort the geometry induced by the metric $g$ - i.e. it might not be \say{metric compatible}. A fundamental result of Riemannian geometry, however, is that for any Riemannian manifold $(\man, g)$ there exists a \textit{unique} connection that is torsion free and is \say{metric compatible}. 

\begin{restatable}[The fundamental theorem of Riemannian Geometry, \cite{petersen2006riemannian}]{thm}{}
    If $\man$ is endowed with a Riemannian metric $g$, then there exists a unique connection $\nabla$ that is
    \begin{enumerate}
        \item \textbf{Torsion free:} $\nabla_XY - \nabla_YX = [X,Y]$, where $[X,Y]$ is the Lie bracket of $X$ and $Y$. 
        \item \textbf{Metric compatible:} $Xg(Y, Z) = g(\nabla_XY, Z) + g(Y, \nabla_XZ).$ \textcolor{red}{gm:what is $g(\cdot,\cdot)$} \tls{the metric tensor, ie the inner product that defines the metric structure of the manifold}
    \end{enumerate}
    This connection is called the Levi-Civita connection.
\end{restatable}

The object $[X, Y]$ is the \textit{Lie bracket} of the vector fields $X$ and $Y$. For any $f \in C^{\infty}(\man)$ it is defined as
\begin{equation}
    [X, Y](f) = X(Y(f)) - Y(X(f))
\end{equation}
and it essentially measures the noncommutativity of the directional derivatives. Concretely, it tells you how much flowing along $X$ then $Y$ differs from flowing along $Y$ then $X$. 

\subsection{A connection on the Wasserstein manifold}


In classical Riemannian geometry, once you have a metric tensor $g$ for a manifold $\man$, the Levi-Civita connection is guaranteed to exist and is uniquely given by something called the \textit{Koszul formula}. Since $\mathcal{P}_2(\mathbb{R}^d)$ is \textit{not} a Riemannian manifold, we do not get the Levi-Civita connection for free. We need to construct a connection and \textit{verify} that the desirable properties (torsion-free and metric compatibility) are satisfied. This construction will follow that put forth by \citet{clancy2021interpolating}. 

Let $D_t$ denote the covariant derivative (which we want to define), and suppose we have a curve $\mu_t$ and two vector fields along it $v_t, w_t$, as well as a tangent field $u_t$. Note that the operator $d/dt$ now corresponds to the directional derivative along $\mu_t$ \textcolor{red}{gm: is $\mu_t$ the the same as $p_t$?}. With this in mind, $D_t$ should satisfy 
\[\frac{d}{dt}\langle v_t, w_t\rangle_{\mu_t} = \langle D_tv_t, w_t\rangle_{\mu_t} + \langle v_t, D_tw_t\rangle_{\mu_t}\]
by metric compatibility. Now we can expand the left-hand side as follows,
\begin{align}
    \frac{d}{dt}\langle v_t, w_t\rangle_{\mu_t} &= \frac{d}{dt}\int_{\mathbb{R}^d}\langle v_t, w_t\rangle d\mu_t \\
    &= \int_{\mathbb{R}^d}\langle \partial_t v_t, w_t\rangle d\mu_t + \int_{\mathbb{R}^d}\langle v_t, \partial _tw_t\rangle d\mu_t + \int_{\mathbb{R}^d}\langle v_t, w_t\rangle d(\partial_t\mu_t).
\end{align}
By the continuity equation, $\partial_t\mu_t = -\nabla \cdot (\rho_tu_t)$ (where $\nabla$ is the Euclidean gradient). Thus,
\begin{align}
    \int_{\mathbb{R}^d}\langle v_t, w_t\rangle d(\partial_t\mu_t) &= -\int_{\mathbb{R}^d}\langle v_t, w_t\rangle d(\nabla \cdot (\rho_tu_t)) \\
    &= \int_{\mathbb{R}^d} \left\langle \nabla \langle v_t, w_t\rangle, u_t\right\rangle  d\mu_t \\
    &= \int_{\mathbb{R}^d} \langle   \nabla v_t \cdot u_t, w_t\rangle d\mu_t   +\int_{\mathbb{R}^d}\langle  \nabla w_t \cdot  u_t , v_t\rangle  d\mu_t 
\end{align}
where the second step follows from the fact that $\int f d(\nabla\cdot U\mu) = \int\langle\nabla f, U\rangle d\mu$ (Florian has a clear proof of why in his notes). Plugging back in yields
\begin{align}
    \frac{d}{dt}\langle v_t, w_t\rangle_{\mu_t} &= \langle\partial_tv_t + \nabla v_t \cdot u_t, w_t\rangle_{\mu_t} + \langle \partial_tw_t + \nabla w_t\cdot u_t, v_t\rangle_{\mu_t}.
\end{align}
By metric compatibility, $D_t$ would therefore have to satisfy
\begin{equation}
    D_tv_t = \partial_tv_t + \nabla v_t \cdot u_t.
\end{equation}
But an important thing to note is that the right-hand side need not be in $T_{\mu_t}\mathcal{P}_2$. Thus, Clancy defines $D_t$ to be the total derivative, and the covariant derivative is
\begin{equation}
    \mathbf{D}_tv_t = \Pi_{\mu_t}\left(D_tv_t\right)
\end{equation}
where $\Pi_{\mu_t}$ is an orthogonal projection onto $T_{\mu_t}\mathcal{P}_2$. To obtain such projection, we use Theorem 8.3.1 of \citet{ambrosio2007gradient} which establishes that for all $v \in L^{2}_{\mu_t}(\mathbb{R}^d)$ the orthogonal projection onto $T_{\mu_t}\mathcal{P}_2$ is characterized by
\begin{equation}
    \nabla \cdot ((v - \Pi_{\mu_t}(v))\mu_t) = 0.
\end{equation}
This should be intuitive -- the component of $v$ orthogonal to the tangent space at $\mu_t$ (given by $v - \Pi_{\mu_t}(v)$) has zero divergence. To compute $\Pi_{\mu_t}(v)$ we simply take a Helmholtz decomposition of its argument $v = \nabla \varphi + w$ and discard the divergence free component $w$. Thus, our covariant derivative of $v_t$ along $u_t$ is given by 
\begin{equation}
    \mathbf{D}_tv_t = \Pi_{\mu_t}\left(\partial_tv_t + \nabla v_t \cdot u_t\right)
\end{equation}
where $\Pi_{\mu_t}$ discards the divergence free component of its argument via a Helmholtz decomposition. I'll note that showing the torsion free property of $\mathbf{D}_t$ is technical, so I'll save that for the HK case only. 

\subsection{A connection on HK by analogy}

Now we'll construct the covariant derivative for the Wasserstein-Fisher-Rao manifold. This too will follow the construction of \citet{clancy2021interpolating}. Let $s_{1, t} = (v_{t}, g_{t})$ and $s_{2, t} = (w_{t}, \alpha_{t})$ be tangent fields along a curve $\rho_t$ with derivative $(u_t, \beta_t)$. Then
\begin{align}
    \frac{d}{dt}\langle \mathbf{s}_{1,t}, \mathbf{s}_{2,t}\rangle_{\rho_t} &= \frac{d}{dt}\int_{\mathbb{R}^d} \left(\langle v_{t}, w_{t} \rangle + 4g_{t}\alpha_{t}\right)\,d\rho_t \\
    &= \int_{\mathbb{R}^d} \frac{d}{dt}\langle v_{t}, w_{t} \rangle \,d\rho_t + \int_{\mathbb{R}^d}\frac{d}{dt}(4g_{t}\alpha_{t})\,d\rho_t \\
    &= \int_{\mathbb{R}^d} \langle \partial_tv_{t}, w_{t} \rangle \,d\rho_t + \int_{\mathbb{R}^d} \langle v_{t}, \partial_tw_{t} \rangle \,d\rho_t + \int_{\mathbb{R}^d} \langle v_{t}, w_{t} \rangle \,d(\partial_t\rho_t) \\
    &\hspace{25mm}+ \int_{\mathbb{R}^d}4\partial_tg_{t}\alpha_{t}\,d\rho_t + \int_{\mathbb{R}^d}4g_{t}\partial_t\alpha_{t}\,d\rho_t + \int_{\mathbb{R}^d}4g_{t}\alpha_{t}\,d(\partial_t\rho_t) \\
    &= \int_{\mathbb{R}^d} \langle \partial_tv_{t}, w_{t} \rangle  + \langle v_{t}, \partial_tw_{t} \rangle  + 4\partial_tg_{t}\alpha_{t} + 4g_{t}\partial_t\alpha_{t}\,d\rho_t + \int_{\mathbb{R}^d}\langle v_t, w_t\rangle +  g_{t}\alpha_{t}\,d(\partial_t\rho_t) \\
    &= T_1 + T_2.
\end{align}
By the continuity equation $\partial_t\rho_t = 4\beta_t\rho_t - \nabla\cdot(\rho_t u_t)$, which implies \textcolor{red}{gm: why there are no $\rho_t d\beta_t$ when going from 28 to 29?} \tls{I think because we are integrating against a measure whose density is $\beta_t\rho_t$. So its not a derivative, and thus product rule does not need to be applied.}
\begin{align}
    T_2 &= \int_{\mathbb{R}^d} \langle v_t, w_t\rangle  + 4g_{t}\alpha_{t}\,d(4\beta_t\rho_t -\nabla\cdot(\rho_tu_t)) \\
    &= \int_{\mathbb{R}^d} 4\langle v_t, w_t\rangle  + 16g_{t}\alpha_{t}\,d(\beta_t\rho_t) - \int_{\mathbb{R}^d} \langle v_t, w_t\rangle  + 4g_{t}\alpha_{t}\,d(\nabla\cdot(\rho_tu_t)) \\
    &= \int_{\mathbb{R}^d} 4\beta_t\langle v_t, w_t\rangle + 16\beta_tg_{t}\alpha_{t} +  \left\langle \nabla \left(\langle v_t, w_t\rangle  + 4g_{t}\alpha_{t}\right), u_t\right\rangle \,d\rho_t \\
    &= \int_{\mathbb{R}^d} 4\beta_t\langle v_t, w_t\rangle + 16\beta_tg_{t}\alpha_{t} +  \left\langle \nabla \langle v_t, w_t\rangle, u_t\right\rangle +  \left\langle \nabla (4g_t\alpha_t), u_t\right\rangle  \,d\rho_t \\
    &= \int_{\mathbb{R}^d} 4\beta_t\langle v_t, w_t\rangle + 16\beta_tg_{t}\alpha_{t} +  \left\langle \nabla v_t \cdot u_t, w_t\right\rangle  + \left\langle\nabla w_t \cdot u_t, v_t\right\rangle + \left\langle\nabla(4g_{t}\alpha_{t}), u_t\right\rangle \,d\rho_t.
\end{align}
Continuing to expand yields
\begin{align}
    T_2 &= \int_{\mathbb{R}^d} 4\beta_t\langle v_t, w_t\rangle + 16\beta_tg_{t}\alpha_{t} +  \left\langle \nabla v_t \cdot u_t, w_t\right\rangle  + \left\langle\nabla w_t \cdot u_t, v_t\right\rangle + 4\alpha_{t}\left\langle\nabla g_{t}, u_t\right\rangle + 4g_{t}\left\langle \nabla \alpha_{t}, u_t\right\rangle \,d\rho_t.\end{align}
Now metric compatibility implies \textcolor{red}{gm:you are not commenting on $T_1$ because there is no ambiguity?} \tls{$T_1$ is already expanded out, so the next lines follow from matching.}
\begin{align}
    \frac{d}{dt}\langle \mathbf{s}_{1,t}, \mathbf{s}_{2,t}\rangle_{\rho_t} &= \langle D_t\mathbf{s}_{1,t}, \mathbf{s}_{2,t}\rangle_{\rho_t} + \langle \mathbf{s}_{1,t}, D_t\mathbf{s}_{2,t}\rangle_{\rho_t} \\
    &= \int_{\mathbb{R}^d}\langle D_t v_t, w_t\rangle + 4(D_tg_t)\alpha_t  + \langle v_t, D_t w_t\rangle + 4g_t(D_t\alpha_t) \, d\rho_t 
\end{align}
Recall from the properties of the solution of the variational problem in \cref{eq: HK tangent vectors} that $\nabla g_t = v_t$ and $\nabla \alpha_t = w_t.$ Unlike in the Wasserstein case, we now have terms that are ambiguous -- they could come from either the first or second term in the metric compatibility condition. In particular, the ambiguous terms are 
\begin{align}
    4\alpha_t\langle \nabla g_t, u_t\rangle \qquad \text{and} \qquad 4g_t\langle \nabla \alpha_t, u_t\rangle. 
\end{align}
One can verify that for any $p \in \mathbb{R}$ the following setting is satisfactory
\begin{equation}
    \mathbf{D}_t\mathbf{s}_t = \Pi_{\rho_t}\left(\begin{bmatrix}
        \partial_tv_t + \nabla v_t\cdot u_t + 2\beta_tv_t +  4pg_tu_t \\
        \partial_tg_t + (1-p)\langle\nabla g_t, u_t\rangle + 2\beta_t g_t
    \end{bmatrix}\right)
\end{equation}
where $\Pi_{\rho_t}$ is an orthogonal projection onto the tangent space at $\rho_t$. \citet{clancy2021interpolating}
shows that the choice of $p = 1/2$ yields a torsion free connection. To define the orthogonal projection we'll again decompose the vectorial portion of the covariant derivative via a Helmholtz decomposition. Since the discarded component of the vector field is divergence free, removing it does not change the mass -- therefore the mass creation term is left unchanged in the projection.  

\subsection{Parallel transport via covariant derivative}

Now we are in a position where we can define parallel transport. Intuitively, a vector field along a curve is a parallel transport if its covariant derivative along the curve is zero.
\begin{restatable}[\citet{lee2018introduction}]{defn}{}
    Let $\man$ be a smooth Riemannian manifold with a connection $\nabla$. A smooth vector field $X$ along a smooth curve $\gamma$ is said to be parallel along $\gamma$ with respect to $\nabla$ if $\mathbf{D}_tX \equiv 0$.
\end{restatable}
We'll first present parallel transport for the Wasserstein case, then follow with the Wasserstein-Fisher-Rao case. Consider a smooth curve $\mu_t$ indexed by $t \in (0, 1)$ with derivative $u_t$ and a vector field $v_t$ along the curve. The vector field $v_t$ is parallel along $\mu_t$ if
\begin{align}
    0 &\equiv \Pi_{\mu_t}\left(\partial_tv_t + \nabla v_t \cdot \nabla\varphi_t\right) \qquad \text{for a.e. }t \in(0,1).
\end{align}
This is tantamount to the requirement that $\text{div}\left(\mu_t\left(\partial_tv_t + \nabla v_t \cdot \nabla \varphi_t\right)\right) \equiv 0$.  Expanding this condition yields
\begin{align}
    \nabla \cdot \left(\mu_t(\partial_tv_t + \nabla v_t \cdot \nabla\varphi_t)\right)  = 0
\end{align}
for almost every $t \in (0,1)$. If we further restrict our attention to tangent vector fields $v_t = \nabla \psi_t$ then we recover the characterization of \citet{ambrosio2008construction} eq. (5.22) 
\begin{equation}
    \nabla \cdot \left((\partial_t\nabla \psi_t + \nabla^2\psi_t\cdot \nabla \varphi_t)\mu_t\right) = 0.
\end{equation}
Now we'll do the same for the HK case. Consider a smooth curve $\rho_t$ indexed by $t \in (0,1)$ with derivative $\mathbf{s}_t = (\nabla\varphi_t, \varphi_t)$. A vector field $(v_t, g_t)$ is parallel along $\rho_t$ if
\begin{align}
    \Pi_{\rho_t}\left(\begin{bmatrix}
        \partial_tv_t + \nabla v_t\cdot \nabla\varphi_t + 2\varphi_tv_t +  2g_t\nabla\varphi_t \\
        \partial_tg_t + \frac{1}{2}\langle\nabla g_t, \nabla\varphi_t\rangle + 2\varphi_t g_t
    \end{bmatrix}\right) = 0
\end{align}
for a.e. $t \in (0,1).$ As mentioned before, the orthogonal projection in the HK case is a Helmholtz decomposition of the vectorial argument (discarding the zero divergence component) and it leaves the mass creation term unchanged. Thus, we are left with the two conditions
\begin{equation}
    \nabla \cdot \left(\rho_t\left(\partial_tv_t + \nabla v_t\cdot \nabla\varphi_t + 2\varphi_tv_t +  2g_t\nabla\varphi_t\right)\right) = 0
\end{equation}
and 
\begin{equation}
    \partial_tg_t + \frac{1}{2}\langle\nabla g_t, \nabla\varphi_t\rangle + 2\varphi_t g_t = 0.
\end{equation}
If we take the vector field $(v_t, g_t)$ to be parallel again, then it can be written as $(\nabla\psi_t, \psi_t)$ and we have
\begin{align}
    \nabla \cdot \left(\rho_t\left(\partial_t\nabla \psi_t + \nabla^2\psi_t \cdot \nabla\varphi_t + 2\varphi_t \nabla \psi_t +  2\psi_t\nabla\varphi_t\right)\right) &= 0 \\
    \text{and} \qquad \partial_t\psi_t + \frac{1}{2}\langle\nabla \psi_t, \nabla\varphi_t\rangle + 2\varphi_t \psi_t &= 0.
\end{align}

\section{The cone metric}
An equivalent construction of the Wasserstein Fisher Rao space comes through lifting measures on the base space $\mathbb{R}^d$ to $\mathfrak{C}(\mathbb{R}^d)$, the metric cone of the base space ($\mathbb{R}^d$). 
\begin{restatable}[Metric cone]{defn}{}
    For a Riemannian manifold $(\man, g)$, the metric cone is
    \[\mathfrak{C}(\man) = \man \times \mathbb{R}_{+} / \{(x, 0) \sim (y, 0)\}\]
    where a point in $C(\man)$ is written as $(x, r)$. The metric tensor on the cone is given by
    \[g_{(x,r)} = dr^2 + r^2g_x.\]
\end{restatable} 
The metric (as one would have in a metric space) is given by
\begin{equation}
    d_{\mathfrak{C}}((x_0, r_0), (x_1, r_1))^2 = r_0^2 + r_1^2 - 2r_0r_1\cos(|x_0 - x_1| \land \pi).
\end{equation}
\tls{Its not clear to me that this metric follows directly the metric tensor. I believe the clipping in the argument of the cosine is an artificial addition...} \ndfg{Check pages 88 and following in Burago Burago Ivanov for exactly this!} Now, given a measure $\lambda \in M(\mathfrak{C})$ we can project it to a measure $\mathfrak{B}\lambda \in M(\mathbb{R}^d)$ with
\begin{equation}
    \int f(x)d\mathfrak{B}\lambda(x) = \int r^2 f(x) d\lambda(x, r) \qquad \text{for all test functions }f.
\end{equation}
Clearly, there are many possible ways to define a lift of a measure -- any $\lambda$ of the form
\begin{equation}
    \lambda(dx, dr) = \eta(x, dr)\mathfrak{B}\lambda(dx) \qquad \text{with } \qquad \int_0^\infty r^2\eta(x, dr) = 1 \text{ \quad for $\mathfrak{B}\lambda$ a.e. $x$}
\end{equation}
will do. To see this plug into the projection formula above and we obtain
\[\int r^2 f(x)\eta(x, dr)\mathfrak{B}\lambda(dx) = \int f(x) \left(\int r^2 \eta(x, dr) \right) \mathfrak{B}\lambda(dx) = \int f(x) \mathfrak{B}\lambda(dx) \]
as desired. \citet{liero2016optimal} show that one can characterize the Wasserstein-Fisher-Rao distance on a base space to the Wasserstein on the lifted cone space,
\begin{equation} \label{eq: cone wasserstein and hk equivalence}
    \text{HK}(\mu_0, \mu_1) = \min\left\{W_{\mathfrak{C}}(\lambda_0, \lambda_1) \,\Big|\, \mathfrak{B}\lambda_0 = \mu_0, \mathfrak{B}\lambda_1 = \mu_1\right\}.
\end{equation}
% \tls{If we want to compute HK distances we need to use the logarithmic entropy transport functional approach. I'll add a section for this soon.}

\subsection{Logarithmic-entropy-transport functional approach}
To compute the HK distance we will use a equivalent characterization that is more computationally amenable. Let $\eta_j = \Pi^j_\#\eta$ and define the Hellinger-Kantorovich entropy-transport functional as
\begin{equation}
    E(\eta; \mu_0, \mu_1) = \int_{\mathbb{R^d}}F\left(\frac{d\eta_0}{d\mu_0}\right)d\mu_0 + \int_{\mathbb{R^d}}F\left(\frac{d\eta_1}{d\mu_1}\right)d\mu_1 + \int_{\mathbb{R}^d\times \mathbb{R}^d}c(|x_0-x_1|)d\eta
\end{equation}
where
\begin{equation}
    c(L) = \begin{cases}
        -2\log(\cos(L)) & L < \pi/2 \\
        \infty & L \geq \pi/2
    \end{cases}
\end{equation}
and 
\begin{equation}
    F(\rho) = \rho\log\rho - \rho + 1.
\end{equation}
\citet{liero2016optimal} show that 
\begin{align} \label{eq: let hk equivalence}
    \text{HK}(\mu_0, \mu_1)^2 &= \inf\left\{E(\eta;\mu_0, \mu_1) \, |\, \eta  \in M(\mathbb{R}^d\times \mathbb{R}^d), \nu_j  \ll \mu_j\right\}.
\end{align}
They also show that the logarithmic-entropy-transport functional gives us a way to compute optimal lifts in the sense of \cref{eq: cone wasserstein and hk equivalence}. In particular, assume that $\eta \in M(\mathbb{R}^d\times \mathbb{R}^d)$ is a minimizer of $E(\cdot\,; \mu_0, \mu_1)$ and for the marginals $\eta_i$ consider the Lebesgue decomposition $\mu_i = \sigma_i\eta_i + \mu_i^\perp$. Then the transport plan $\gamma_{\eta} \in M(\mathfrak{C}_{\mathbb{R}^d} \times \mathfrak{C}_{\mathbb{R}^d})$ defined by
\begin{multline}
    \gamma_{\eta}(dz_0, dz_1) = \delta_{\sqrt{\sigma_0(x_0)}}(dr_0)\delta_{\sqrt{\sigma_1(x_1)}}(dr_1)\eta(dx_0, dx_1) \\
    + \delta_1(dr_0)\mu_0^\perp(dx_0)\delta_{\mathfrak{o}}(dz_1) + \delta_{\mathfrak{o}}(dz_0)\delta_1(dr_1)\mu_1^\perp(dx_1) 
\end{multline}
and the associated lifts $\lambda_i = \Pi_{\#}^i\gamma_\eta$ are optimal for \cref{eq: cone wasserstein and hk equivalence}. This means we can solve the logarithmic-entropy-transport functional approach and use the optimal plan to construct appropriate lifts of the original measures onto the cone. 

\subsection{Relevant geometric objects on the cone}

Taking $\man = \mathbb{R}^d$ to be the base space of our cone construction, we can construct tangent vectors as
\begin{equation}
    (u, a) \in T_{(x,r)}\mathfrak{C}(\mathbb{R}^d)
\end{equation}
where $u \in T_x\mathbb{R}^d \cong \mathbb{R}^d$ and $a \in \mathbb{R}$. Equivalently, it can be written as $u + a\partial_r$ where $\partial_r$ is the canonical radial vector field. The metric cone described above falls under the \say{warped product} construction, which we define explicitly below. 
\begin{restatable}[Warped product, \citet{petersen2006riemannian}]{defn}{}
    Given a Riemannian metric $(\man, g)$, a warped product (over $I$) is defined as a metric on $I \times \man$, where $I \subset \mathbb{R}$ is an open interval with metric 
    \[g = dr^2 + \rho^2(r)g\]
    where $\rho > 0$ on $I$. This is sometimes denoted $I \times_{\rho} \man$, where $I$ is called the base space and $\man$ is called the fiber.
\end{restatable}
Having established this definition, it is clear that the cone construction is a warped product using $\rho(r) = r$. With this in mind, we can used the theory of warped products to study the cone metric. Before moving on, we need to establish the \say{lift} of a vector field on a manifold to a vector field on a product manifold.

\begin{restatable}[Lifting, \citet{o1983semi}]{defn}{}
    Let $\pi: \mathcal{N} \times \man\rightarrow \mathcal{N}$ be the projection onto the first factor $\pi(p, q) = p$. 
    \begin{enumerate}
        \item If $x \in T_p\mathcal{N}$ and $q \in \mathcal{M}$ then the lift $\tilde{x}$ of $x$ to $(p,q)$ is the unique vector in $T_{(p,q)}(\mathcal{N} \times q)$ such that $d\pi(\tilde{x}) = x$. We denote the set of all such horizontal lifts $\mathfrak{L}(\mathcal{N}).$
        \item If $A \in \mathfrak{X}(\mathcal{N})$, then the lift of $A$ to $\mathcal{N} \times \man$ is the vector field $B$ whose values at each $(p,q)$ is the lift of $A_p$ to $(p,q).$
    \end{enumerate}
    Vertical lifts are defined in the same way but using the projection onto the second factor $\sigma: \mathcal{N} \times \man \rightarrow \man$. The set of all vertical lifts are denoted $\mathfrak{L}(\man).$ 
\end{restatable}



\subsubsection{The Levi-Civita connection}
Proposition 35 of \citet{o1983semi} derives the Levi-Civita connection for warped products. 

\begin{restatable}[Proposition 35, \citet{o1983semi}]{prop}{}
    Let $\man$ and $I$ be Riemannian manifolds with connections denoted $^{\man}\nabla$ and $^I\nabla$ respectively. On $I \times_\rho \man$, if $X, Y \in \mathfrak{L}(\man)$ and $V, W \in \mathfrak{L}(I)$ then
    \begin{enumerate}
        \item $\nabla_VW \in \mathfrak{L}(I)$ is the horizontal lift of $^{I}\nabla_VW$ on $I$.
        \item $\nabla_VX = \nabla_XV = \left(\frac{V(\rho)}{\rho}\right)X$.
        \item $\nabla_XY = \,^\man\nabla_XY - \left(\frac{\langle X, Y\rangle}{\rho}\right)\text{grad}(\rho)$.
    \end{enumerate}
\end{restatable}
\noindent Now we can directly apply this proposition to compute the connection on $\mathfrak{C}(\mathbb{R}^d).$ 

\begin{restatable}[Levi-Civita connection on $\mathfrak{C}(\mathbb{R}^d)$]{corollary}{} \label{corollary: connection on cone}
    Let $X, Y \in \mathfrak{L}(\mathbb{R}^d)$ and let $V, W \in \mathfrak{L}(\mathbb{R}_+)$. In particular, let $X = (X^1(x), \dots, X^n(x))$ and $Y= (Y^1(x), \dots, Y^d(x))$ in standard coordinates. Let $\nabla$ denote the Levi-Civita connection on $\mathfrak{C}(\mathbb{R}^d) \setminus \{\mathfrak{o}\}$ where $\mathfrak{o}$ is the base of the cone and let $(x,r) \in \mathfrak{C}(\mathbb{R}^d)$. It holds that 
    \begin{enumerate}
        \item $\nabla_VW$ is given by 
        \[\nabla_VW(x,r) = \left(0, V(r)\frac{\partial W(r)}{\partial r}\right).\]
        \item $\nabla_VX(x,r) = \nabla_XV(x,r) = \frac{1}{r}X(x,r)$.
        \item Finally,
        \[\nabla_XY(x,r) = \left(\sum_{i = 1}^dX^i(x) \frac{\partial Y(x)}{\partial x^i},  - \frac{\langle X(x), Y(x) \rangle_{\mathfrak{C}}}{r}\right). \]
    \end{enumerate}
\end{restatable}

\subsubsection{Geodesics}

Fortunately, geodesics on the cone space are available in closed form, and we will make them explicit here \cite{liero2016optimal}. Given two points $z_0 = (x_0, r_0)$ and $z_1 = (x_1, r_1)$, the geodesic between them can be expressed as
\begin{equation} \label{eq: cone geodesics}
    z(s) = \left[x(s; z_0, z_1) , r(s; z_0, z_1)\right]
\end{equation}
where
\begin{equation} \label{eq: cone geodesic radial}
    r(s; z_0, z_1)^2 = (1-s)^2r_0^2 + s^2r_1^2 + 2s(1-s)r_0r_1\cos(\min\{\pi, \|x_1 - x_0\|_{2}\})
\end{equation}
and 
\begin{equation} \label{eq: cone geodesic base}
    x(s; z_0, z_1) = (1 - \rho(s; z_0, z_1))x_0 + \rho(s; z_0, z_1)x_1
\end{equation}
with
\begin{equation}
    \rho(s; z_0, z_1) = \begin{cases}
        \frac{1}{\|x_1 - x_0\|_2}\arccos\left(\frac{(1-s)r_0 + sr_1\cos|x_1 - x_0|}{r(s; z_0, z_1)}\right) & \text{for }\|x_1 - x_0\|_2 < \pi \\
        \frac{1}{2}\left(1 + \text{sign}((1-s)r_0 + sr_1)\right) & \text{otherwise}
    \end{cases}.
\end{equation}
Now, given an optimal coupling $\gamma$ between two (probability) measures $\lambda_1, \lambda_2 \in M(\mathfrak{C})$, we can compute the Wasserstein geodesic interpolant between them as $z(s; \cdot, \cdot)_{\#}\gamma$. One can also obtain these equations through the Levi-Civita connection defined in the previous section, as it holds that for any geodesic $\gamma$, $\nabla_{\dot \gamma}\dot \gamma = 0$. Expanding this equation using the decomposition $\dot \gamma = v + \dot r\partial_r$ \Cref{corollary: connection on cone} yields
\begin{align}
    0 = \nabla_{\dot \gamma}\dot \gamma &= \nabla_{v + \dot r \partial_r}(v + \dot r\partial_r) \\
    &= \nabla_vv + \nabla_v \dot r \partial_r + \dot r \nabla_{\partial_r}v + \dot r \nabla_{\partial_r}(\dot r \partial_r) \\
    &= \dot v - r\|v\|_2^2\partial_r + \ddot r \partial_r + \dot r \nabla_v\partial_r + \dot r \nabla_{\partial_r}v + 0\cdot  \partial_r + \dot r^2 \nabla_{\partial_r}\partial_r \\
    &= \left(\dot v + 2v\frac{\dot r}{r}\right) + \left(\ddot r  - r\|v\|_2^2\right)\partial_r.
\end{align}
Thus, the geodesic equations on the cone are
\begin{equation}\label{eq: geodesic equations on the cone}
    \dot v + 2v\frac{\dot r}{r} = 0 \qquad \text{and} \qquad \ddot r  - r\|v\|_2^2 = 0.
\end{equation}
This also has an interesting consequence. If we differentiate the following quantity against $t$
\begin{align*}
    \frac{d}{dt}(r^2v) &= 2r\dot rv + r^2 \dot v \\
    &= 2r \dot r v + r^2(-2v\dot r/r) \\
    &= 0.
\end{align*}
Thus $r^2v$ is constant in time. We will leverage this property to simplify the Jacobi equation. 

\subsubsection{The curvature tensor}

The last piece of the puzzle is the Riemann curvature tensor, as we want to be able to compute $R(u, \dot\gamma(t))\dot\gamma(t)$ to solve the Jacobi equation described in the next subsection. The curvature tensor of a (pseudo)-Riemannian manifold measures the noncommutativity of the covariant derivative, and is defined as
\begin{equation}
    R(X,Y)Z = \nabla_X \nabla_YZ - \nabla_Y\nabla_XZ - \nabla_{[X,Y]}Z
\end{equation}
where $X,Y, Z \in \mathfrak{X}(\man)$ are vector fields over the manifold. In the situation of interest, we have $\man = \mathfrak{C}(\mathbb{R}^d)$ - we will appeal to another result from the theory of warped-product manifolds to obtain an expression for the Riemann curvature tensor as a function of the curvature tensors of the base space and the fiber. 
\begin{restatable}[Proposition 7.42, \citet{o1983semi}]{prop}{} \label{prop: warped product curvature tensor}
      Let $\man$ and $I$ be Riemannian manifolds with Riemann curvature tensors denoted $^{\man}R$ and $^IR$ respectively. On $I \times_\rho \man$, if $X, Y, Z \in \mathfrak{L}(\man)$ and $U, V, W \in \mathfrak{L}(I)$ then
    \begin{enumerate}
        \item $R(U, V)W \in \mathfrak{L}(I)$ is the horizontal lift of $^IR(U,V)W$ on $I$.
        \item $R(Y,U)V = (H^\rho(U,V)/\rho)Y$ where $H^\rho$ is the Hessian of $\rho$.
        \item $R(U,V)Y = R(Y,Z)U = 0.$
        \item $R(U,Y)Z = (\langle Y,Z\rangle/\rho)\nabla_U\,\text{grad}(\rho)$\item $R(Y,Z)X = \,^\man R(Y,Z)X - (\langle \text{grad}\rho, \text{grad}\rho\rangle/\rho^2)\left\{\langle Y,X \rangle  Z - \langle Z, X\rangle Y\right\}.$
    \end{enumerate}
\end{restatable}
Now we'll drop the dependence on time and write $J = j +  b\partial_r$ and $\dot\gamma = v +  \dot r\partial_r$ where $j, v \in \mathfrak{L}(\mathbb{R}^d)$ are vertical lifts, and $b\partial_r,\dot r\partial_r \in \mathfrak{L}(I)$ are horizontal lifts. By linearity
\begin{multline}
    R(J, \dot\gamma)\dot\gamma = R(j, v)v + \dot rR(j, \partial_r)v + bR(\partial_r, v)v + b\dot rR(\partial_r, \partial_r)v \\
    + \dot rR(j, v)\partial_r + \dot r^2R(j, \partial_r)\partial_r + \dot rbR(\partial_r, v)\partial_r + b\dot r^2R(\partial_r, \partial_r)\partial_r.
\end{multline}
\Cref{prop: warped product curvature tensor} indicates that any term $R(a,b)c$ where $a,b \in \mathfrak{L}(\mathbb{R}^d)$ and $c \in \mathfrak{L}(I)$ (or vice versa) is zero. Also note that $I$ is flat, so $^IR \equiv0$, implying $r^2R(\partial_r, \partial_r)\partial_r = 0$. Thus,
\begin{align}
    R(J, \dot\gamma)\dot\gamma = R(j, v)v + \dot rR(j, \partial_r)v + bR(\partial_r, v)v + \dot r^2R(j, \partial_r)\partial_r + \dot rbR(\partial_r, v)\partial_r.
\end{align}
\Cref{prop: warped product curvature tensor} indicates that 
\begin{align}
    \dot rR(j, \partial_r)v &= -\dot rR(\partial_r, j)v & \text{by antisymmetry}\\
    &= \left(\frac{\langle j, v\rangle_{\mathbb{R}^d}}{\rho}\right)\nabla_{\partial_r}\text{grad}(\rho) \\
    &= \left(\frac{\langle j, v\rangle_{\mathbb{R}^d}}{\rho}\right)\nabla_{\partial_r}\partial_r \\
    &= 0
\end{align}
and $bR(\partial_r, v)v = 0$ since $\rho(r) = r$. We also have
\begin{align}
    \dot r^2R(j, \partial_r)\partial_r &= \left(\frac{H^{\rho}(\partial_r, \partial_r)}{\rho}\right)j \\
    &= 0
\end{align}
since $H^{\rho}\equiv 0$ and $\dot rbR(\partial_r, v)\partial_r = 0$ by the same logic. Putting it together, we are left with
\begin{align}
    R(J, \dot\gamma)\dot\gamma &= R(j,v)v \\
     &= \,^{\mathbb{R}^d}R(j,v)v - \frac{\langle \text{grad} \rho, \text{grad}\rho\rangle}{\rho^2}\left(\langle j, v \rangle v -  \langle v, v\rangle j\right) \\
     &= -\frac{1}{r^2}\left(\langle j, v\rangle_{\mathfrak{C}} \cdot v - \|v\|^2_{\mathfrak{C}}\cdot j\right) \\
     &= j\|v\|^2_{2} - v\langle j,v\rangle_{\mathbb{R}^d}. \label{eq: jacobi riemann expression}
\end{align}
Now we can plug this directly into the Jacobi equation. Note that this scalar quantity above should really be lifted to the cone (and we will do so in the subsequent derivations).

\subsubsection{Putting it together: Jacobi fields}

To compute Jacobi fields on the cone space described above given initial conditions $J(0)$ and $J'(0)$, one needs to solve the Jacobi equation
\begin{equation}
    \nabla_{\dot\gamma}^2J(t) + R(J(t), \dot{\gamma}(t))\dot{\gamma}(t) = 0
\end{equation}
where $\gamma$ is the geodesic of interest on the cone space and $R$ is the Riemann curvature tensor. Let $J = j + b\partial_r$ and let $\dot\gamma = v + \dot r \partial_r$. Plugging in the result from \cref{eq: jacobi riemann expression} yields
\begin{align}
    \nabla_{v + \dot r \partial_r}^2(j + b\partial_r) + j\|v\|^2_{2} - v\langle j,v\rangle_{\mathbb{R}^d} = 0.
\end{align}
Now we'll expand the first term as follows
\begin{align}
    \nabla_{v + \dot r \partial_r}(j + b\partial_r) &= \nabla_vj + \nabla_v(b\partial_r) + \nabla_{\dot r\partial_r}j + \nabla_{\dot r\partial_r}(b\partial_r) \\
    &= \nabla_vj + v(b)\partial_r + b\nabla_v\partial_r + \dot r\nabla_{\partial_r}j + \nabla_{\dot r\partial_r}\partial_r + 0.
\end{align}
Continuing, we have
\begin{align}
    \nabla_{v + \dot r \partial_r}(j + b\partial_r)&= \,^{\mathbb{R}^d}\nabla_vj - \frac{\langle v, j\rangle_{\mathfrak{C}}}{\rho}\text{grad}\rho + \dot b\partial_r +  \left(\frac{\partial_r(\rho)}{\rho}\right)bv  + \frac{\dot r\partial_r(\rho)}{\rho}j & \text{ by \Cref{corollary: connection on cone}}\\
    &= \left(\dot j + \frac{v}{r}b  + \frac{\dot r}{r}j, \dot b - \frac{\langle v, j\rangle_{\mathfrak{C}}}{r}\right) \\
    &= \left(\dot j + \frac{v}{r}b  + \frac{\dot r}{r}j, \dot b - r\langle v, j\rangle_{\mathbb{R}^d}\right)
\end{align}
For convenience, let $A = \dot j + \frac{v}{r}b  + \frac{\dot r}{r}j$ and let $c = \dot b - r\langle v, j\rangle_{\mathbb{R}^d}$. Now we'll apply the covariant derivative again,
\begin{align}
    \nabla^2_{\dot\gamma}J(t) &= \nabla_{\dot\gamma}(A + c\partial_r) \\
    &= \nabla_v A + \dot r\nabla_{\partial_r}A + \dot c\partial_r + c\nabla_{\dot\gamma}\partial_r\\
    &= \dot A  - \frac{\langle v, A\rangle_{\mathfrak{C}}}{r}\partial_r + \dot r\nabla_{\partial_r}A + \dot c\partial_r + c\nabla_{v}\partial_r + \dot rc\nabla_{\partial_r}\partial_r \\
    &= \dot A  - r\langle v, A\rangle_{\mathbb{R}^d}\partial_r + \frac{\dot r}{r}A + \dot c\partial_r + \frac{cv}{r} \\
    &= \left(\dot A  + \frac{\dot r}{r}A + \frac{c}{r}v\right) + \left(\dot c - r\langle v, A \rangle_{\mathbb{R^d}}\right)\partial_r.
\end{align}
Plugging back into the Jacobi equation yields
\begin{equation}
    \left(\dot A  + \frac{\dot r}{r}A + \frac{c}{r}v + j\|v\|_2^2 - v\langle j, v\rangle_{\mathbb{R^d}}\right) + \left(\dot c - r\langle v, A \rangle_{\mathbb{R^d}}\right)\partial_r = 0.
\end{equation}
Expanding further, we have
\begin{align}
    \dot A &= \ddot j + \frac{r\dot v - v \dot r}{r^2}b + \frac{v}{r}\dot b + j\left(\frac{r\ddot r - \dot r ^2}{r^2}\right) + \frac{\dot r}{r}\dot j\\
    &= \ddot j + \left(\frac{\ddot r}{r} - \frac{\dot r^2}{r^2}\right)j + \frac{\dot r}{r}\dot j + \left(\frac{\dot b}{r} - \frac{b \dot r}{r^2}\right)v + \frac{b}{r}\dot v
\end{align}
and 
\begin{align}
    \dot c &= \ddot b - \dot r \langle v, j\rangle - r\langle \dot v, j \rangle - r\langle v, \dot j\rangle.
\end{align}
The geodesic equations in \cref{eq: geodesic equations on the cone} indicate that $\dot v = -2v\dot r/r$ and $\ddot r = r\|v\|_2^2$. Substituting this in yields
\begin{align*}
    \dot A = \ddot j + \left(\|v\|_2^2 - \frac{\dot r^2}{r^2}\right)j + \frac{\dot r}{r}\dot j + \left(\frac{\dot b}{r} - \frac{b \dot r}{r^2}\right)v  -2v\frac{b\dot r}{r^2}.
\end{align*}
Adding the $A\dot r/r$ term from the Jacobi equation yields
\begin{align*}
    \dot A + \frac{\dot r}{r}A = \ddot j +  2\frac{\dot r}{r}\dot j +\|v\|_2^2j + \left(\frac{\dot b}{r} - \frac{2b \dot r}{r^2}\right)v
\end{align*}
and plugging back into the base component of the Jacobi equation
\begin{align}
    0 &= \ddot j +  2\frac{\dot r}{r}\dot j +\|v\|_2^2j + \left(\frac{\dot b}{r} - \frac{2b \dot r}{r^2}\right)v + \frac{c}{r}v + j\|v\|_2^2 - v\langle j, v\rangle \\
    &= \ddot j +  2\frac{\dot r}{r}\dot j +\|v\|_2^2j + \left(\frac{\dot b}{r} - \frac{2b \dot r}{r^2}\right)v + \left(\frac{\dot b}{r} - \langle v, j\rangle \right)v + j\|v\|_2^2 - v\langle j, v\rangle\\
    &= \ddot j +  2\frac{\dot r}{r}\dot j +2\|v\|_2^2j + \left(\frac{2\dot b}{r} - \frac{2b \dot r}{r^2}\right)v - 2v\langle j, v\rangle \\
    &= \ddot j +  2\frac{\dot r}{r}\dot j +2\|v\|_2^2j + \frac{2}{r}\left(\dot b - b\frac{ \dot r}{r}\right)v - 2v\langle j, v\rangle.
\end{align}
Now we'll do the same for the radial component,
\begin{align*}
    \dot c - r\langle v, A\rangle &= \dot c - r\langle v, \dot j\rangle - b\|v\|_2^2 - \dot r\langle v, j\rangle \\
    &= \ddot b - \dot r \langle v, j\rangle - r\langle \dot v, j \rangle - r\langle v, \dot j\rangle  - r\langle v, \dot j\rangle - b\|v\|_2^2 - \dot r\langle v, j\rangle \\
    &= \ddot b - 2r\langle v, \dot j\rangle - b \|v\|_2^2.
\end{align*}
Thus, our Jacobi equations are

\begin{equation}
    \boxed{\ddot j +  2\frac{\dot r}{r}\dot j +2\|v\|_2^2j + \frac{2}{r}\left(\dot b - b\frac{ \dot r}{r}\right)v - 2v\langle j, v\rangle = 0}
\end{equation}
and 
\begin{equation}
    \boxed{\ddot b - 2r\langle v, \dot j\rangle - b \|v\|_2^2 = 0.}
\end{equation}


\subsubsection{Solving the Jacobi equations}
Since $r^2v$ is constant in time we will define it to be $q := r^2v$ and substitute in $v = q/r^2$, 
\begin{equation}
    \ddot j +  2\frac{\dot r}{r}\dot j +\frac{2}{r^4}\|q\|_2^2j + \frac{2}{r^3}\left(\dot b - b\frac{ \dot r}{r}\right)q - \frac{2}{r^4}q\langle j, q\rangle = 0
\end{equation}
and 
\begin{equation}
    \ddot b - \frac{2}{r}\langle q, \dot j\rangle - \frac{b}{r^4} \|q\|_2^2 = 0.
\end{equation}
Now we will decompose $j$ into $j_{\perp} + aq$, where $j_{\perp} \perp q$. It follows that $\langle j, q\rangle = a\|q\|_2^2$, $\dot j = \dot j_{\perp} + \dot aq$, and $\langle \dot j, q\rangle = \dot a \|q\|_2^2$. Thus,
\begin{equation}
    \ddot j +  2\frac{\dot r}{r}\dot j +\frac{2}{r^4}\|q\|_2^2j_{\perp} + \frac{2}{r^3}\left(\dot b - b\frac{ \dot r}{r}\right)q = 0
\end{equation}
and 
\begin{equation}
    \ddot b - \frac{2\dot a}{r}\|q\|_2^2 - \frac{b}{r^4} \|q\|_2^2 = 0.
\end{equation}
Since the LHS of the equations are equal to zero, we can project onto the orthogonal complement of $q$ and the span of $q$ to get an equivalent form for the first equation,
\begin{equation}
    \ddot j_{\perp} +  2\frac{\dot r}{r}\dot j_{\perp} +\frac{2}{r^4}\|q\|_2^2j_{\perp} = 0
\end{equation}
and 
\begin{equation}
    \ddot a + 2\frac{\dot r}{r}\dot a + \frac{2}{r^3}\left(\dot b - b\frac{ \dot r}{r}\right) = 0.
\end{equation}
Now let $u_\perp = rj_\perp$ and let $u = b/r$. By \Cref{lemma: orthog jacobi equation} we have
\begin{equation} 
    \ddot u_{\perp} + \frac{\|q\|_2^2}{r^4}u_\perp = 0.
\end{equation}
This gives us a way to solve for $u_\perp$. By \Cref{lemma: parallel jacobi equation} we have
\begin{equation} \label{eq: p and w ode}
    \dot p + 2\dot u = 0 \qquad \text{and} \qquad \dot w - 2\frac{\|q\|_2^2}{r^2}p = 0
\end{equation}
where $p = r^2\dot a$ and $w = r^2\dot u$. To solve the equations in \cref{eq: p and w ode} observe that $\dot p = (-2w/r^2)$ and $\dot w = 2p\|q\|_2^2/r^2$. Thus we have the system
\begin{align*}
    \begin{bmatrix}
        \dot p \\ \dot w
    \end{bmatrix} &= 
    \underbrace{
    \begin{bmatrix}
        0 & \frac{-2}{r(t)^2} \\
        2\frac{\|q\|_2^2}{r(t)^2} & 0 
    \end{bmatrix}}_{A(t)}
    \begin{bmatrix}
        p \\
        w
    \end{bmatrix} \\
    &= \frac{1}{r(t)^2}
    \begin{bmatrix}
        0 & -2  \\
        2\|q\|_2^2 & 0 
    \end{bmatrix}
    \begin{bmatrix}
        p \\
        w
    \end{bmatrix}.
\end{align*}
Now we'll define $\tau(t) := \int_0^tds/r^2(s)$, which means 
\[\frac{d}{dt} = \frac{1}{r(t)^2}\frac{d}{d\tau}.\]
Under this change of variables, we obtain
\begin{align*}
    \frac{d}{d\tau}
    \begin{bmatrix}
        p \\ w
    \end{bmatrix} 
    &= \underbrace{\begin{bmatrix}
        0 & -2  \\
        2\|q\|_2^2 & 0 
    \end{bmatrix}}_{M}
    \begin{bmatrix}
        p \\
        w
    \end{bmatrix}.
\end{align*}
The eigenvalues of $M$ are $\lambda = \pm 2ci$ with eigenvectors $\begin{bmatrix}
    \pm i/c & 1
\end{bmatrix}^T$ with $c = \|q\|_2$. Thus the general solution is of the form 
\begin{align*}
    \begin{bmatrix}
        p \\ w
    \end{bmatrix} &= k_1\begin{bmatrix}
        -\frac{1}{c}\sin(2c\tau) \\ \cos(2c\tau)
    \end{bmatrix} + k_2\begin{bmatrix}
        \frac{1}{c}\cos(2c\tau) \\ \sin(2c\tau)
    \end{bmatrix}
\end{align*}
Solving for initial conditions $p_0$ and $w_0$ yields
\begin{align*}
    p(t) &= -\frac{w_0}{\|q\|_2}\sin(2\|q\|_2\tau(t)) + p_0\cos(2\|q\|_2\tau(t)) \\
    w(t) &= w_0\cos(2\|q\|_2\tau(t)) + p_0\|q\|_2\sin(2\|q\|_2\tau(t)).
\end{align*}
Plugging in and integrating yields
\begin{align*}
    a(t) &= a_0 -\frac{w_0}{\|q\|_2} \int_0^t\frac{\sin(2\|q\|_2\tau(s))}{r(s)^2}ds + p_0 \int_0^t \frac{\cos(2\|q\|_2\tau(s))}{r(s)^2} ds \\
    &= a_0 + \frac{w_0}{2\|q\|_2^2}(\cos(2\|q\|_2\tau(t)) - 1) + \frac{p_0}{2\|q\|_2^2}\sin(2\|q\|_2\tau(t)).
\end{align*}
Also note that we have $p_0 = r_0^2\dot a_0$ and $w_0 = r_0^2\dot u_0 = r_0^2 (\frac{\dot b_0}{r_0} - b_0\frac{\dot r_0}{r_0^2}) = \dot b_0r_0 - b_0\dot r_0.$ Putting it together, we have
\begin{equation}
    \boxed{a(t) = a_0 + \frac{\dot b_0r_0 - b_0\dot r_0}{2\|q\|_2^2}(\cos(2\|q\|_2\tau(t))-1) + \frac{r_0^2\dot a_0}{2\|q\|_2^2}\sin(2\|q\|_2\tau(t)).}
\end{equation}
Doing the same for $u$ yields
\begin{align*}
    u(t) &= u_0 + w_0 \int_0^t\frac{\cos(2\|q\|_2\tau(s))}{r(s)^2}ds + p_0 \|q\|_2 \int_0^t \frac{\sin(2\|q\|_2\tau(s))}{r(s)^2} ds \\
    &= u_0 + \frac{w_0}{2\|q\|_2}\sin(2\|q\|_2\tau(s)) - \frac{p_0}{2} (\cos(2\|q\|_2\tau(s))-1).
\end{align*}
This implies
\begin{equation}
    \boxed{b(t) = r(t)\left[\frac{b_0}{r_0} + \frac{\dot b_0r_0 - b_0\dot r_0}{2\|q\|_2}\sin(2\|q\|_2\tau(t)) + \frac{r_0^2\dot a_0}{2} (1- \cos(2\|q\|_2\tau(t)))\right].}
\end{equation}
\tls{one can solve for $a(0)$ and $\dot a(0)$ from ivs of the Jacobi field. Will describe in detail later.} All that remains is to solve for $u_\perp$. 
Fortunately, $\ddot u_{\perp} + \frac{\|q\|_2^2}{r^4}u_\perp = 0$ falls under the umbrella of a well studied variety of ODEs called time-varying harmonic oscillator. In this case, because $r$ is not a constant function of time, this is a harmonic oscillator with varying frequency. The solution is given by \Cref{lemma: time varying oscillator}, and has the form
\begin{equation}
    u_{\perp}(t) = r(t)\left[\frac{u_{\perp,0}}{r_0}\cos(\|q\|_2\sqrt{2}\tau(t)) + \frac{\dot u_{\perp,0} r_0 - u_{\perp,0} \dot r_0}{\|q\|_2\sqrt{2}}\sin(\|q\|_2\sqrt{2}\tau(t))\right].
\end{equation}
which means 
\begin{equation}
    \boxed{j_{\perp}(t) = \frac{r_0j_{\perp,0}}{r_0}\cos(\|q\|_2\sqrt{2}\tau(t)) + \frac{r_0^2\dot j_{\perp, 0}}{\|q\|_2\sqrt{2}}\sin(\|q\|_2\sqrt{2}\tau(t)).}
\end{equation}


\begin{algorithm}[H]
    \begin{algorithmic}[1]
        \caption{Jacobi field on $\mathfrak{C}(\mathbb{R}^d)$}
        \label{alg: jacobi field}
        \Require Initial conditions $J(0) = j_0 + b_0\partial_r$ and $J'(0) = \dot j_0 + \dot b_0 \partial_r$, cone geodesic $\gamma(s) = x(s) + r(s)\partial_r$, terminal time $t$.\\
        Compute geodesic invariant $q_\gamma = \dot x(0)r(0)^2$ and pseudotime $\tau(t) = \int_0^tds/r^2(s)$. \\
        Decompose $j_0 = j_{\perp, 0} + a_0q_{\gamma}$ where $j_{\perp, 0} \perp q_{\gamma}$ and $\dot j_0 = \dot j_{\perp, 0} + \dot a_0q_{\gamma}$ where $\dot j_{\perp, 0} \perp q_{\gamma}$.\\
        Compute 
        \begin{align*}
        j_{\perp}(t) &= \frac{r_0j_{\perp,0}}{r_0}\cos(\|q_{\gamma}\|_2\sqrt{2}\tau(t)) + \frac{r_0^2\dot j_{\perp, 0}}{\|q_{\gamma}\|_2\sqrt{2}}\sin(\|q_{\gamma}\|_2\sqrt{2}\tau(t)) \\
        b(t) &= r(t)\left[\frac{b_0}{r_0} + \frac{\dot b_0r_0 - b_0\dot r_0}{2\|q_{\gamma}\|_2}\sin(2\|q_{\gamma}\|_2\tau(t)) + \frac{r_0^2\dot a_0}{2} (1- \cos(2\|q_{\gamma}\|_2\tau(t)))\right] \\
        a(t) &= a_0 + \frac{\dot b_0r_0 - b_0\dot r_0}{2\|q_{\gamma}\|_2^2}(\cos(2\|q_{\gamma}\|_2\tau(t))-1) + \frac{r_0^2\dot a_0}{2\|q_{\gamma}\|_2^2}\sin(2\|q_{\gamma}\|_2\tau(t)).
        \end{align*}
        \\ \Return $J(t) = j_{\perp}(t) + a(t)q_{\gamma} + b(t)\partial_r$.
    \end{algorithmic}
\end{algorithm}

\section{HK parallel transport via Jacobi fields}

\subsection{Jacobi Fields}
Jacobi fields allow us to measure the effect of curvature on a one parameter family of geodesics. A \textit{Jacobi field} is a vector field along a geodesic $\gamma$ capturing the difference between a geodesic and a perturbed geodesic along the parameter of the family. Concretely, let $\gamma_{\tau}$ be a smooth, one-parameter family of geodesics with $\gamma_0 = \gamma$. Then
\begin{equation}
    J(t) = \frac{\partial}{\partial \tau}\gamma_{\tau}(t)\bigg|_{\tau = 0}
\end{equation}
is a Jacobi field, and it describes the infinitesimal behavior of the geodesic family about $\tau = 0$. Jacobi fields satisfy the so-called Jacobi equation
\begin{equation}
    \mathbf{D}_t^2J(t) + R(J(t), \dot{\gamma}(t))\dot{\gamma}(t) = 0
\end{equation}
where $\mathbf{D}_t$ is the covariant derivative (typically with respect to the Levi-Civita connection) $\dot{\gamma}(t) = d\gamma(t)/dt$, and $R(\cdot, \cdot)$ is the Riemann curvature tensor. A key fact about Jacobi fields is the following, which arises from the properties of the second-order ODE above. 

\begin{restatable}[Uniqueness of Jacobi Fields]{prop}{}
The value of $J(t)$ and $\mathbf{D}_tJ(t)$ at a single point along $\gamma$ completely and uniquely determines the Jacobi field along the entire geodesic.    
\end{restatable}

Observe that this result indicates that there are two ways for one to define Jacobi fields on a manifold. To establish the connection to parallel transport, we will construct a specific Jacobi field as is done by \citet{louis2018fanning}. Consider a geodesic $\gamma(t)$ connecting a point $p = \gamma(0)$ and some $q$, and let $w \in T_p\man$. Now define the $1$-parameter family of geodesics 
\[\exp_{\gamma(t)}(s(\dot{\gamma}(t) + \epsilon w))\]
parameterized by $\epsilon$. As the name \say{fanning} suggests, this family represents a collection of geodesics fanning out from the point $\gamma(t)$ -- a larger $\epsilon$ means the initial velocity is pulled further towards $w$ rather than the original geodesic velocity $\dot{\gamma}(t).$ This parameterized family of geodesics then induces the following Jacobi field,
\begin{align}
    J^{w}_{\gamma(t)}(s) = \frac{\partial}{\partial\epsilon}\bigg|_{\epsilon = 0}\exp_{\gamma(t)}(s(\dot{\gamma}(t) + \epsilon w)) \in T_{\gamma(t+s)}\man.
\end{align}
As shown by \citet{louis2018fanning} it holds that 
\begin{equation}
    P_{0, s}(w) = \frac{J^w_{\gamma(0)}(s)}{s} + O(s^2)
\end{equation}
where $P_{0, s}$ is the vector field obtained by parallel transporting $w$ along the geodesic $\gamma$. In practice, if one can compute the exponential and log maps, then the Jacobi field can be approximated by finite differences,
\[J^w_{\gamma(t)}(s) \cong \frac{\exp_{\gamma(t)}(s(\dot{\gamma}(t) + \epsilon w)) - \exp_{\gamma(t)}(s\dot{\gamma}(t))}{\epsilon}.\]
Unfortunately this construction is not directly applicable to our scenario because neither $\mathcal{P}_2(\mathbb{R}^d)$ or $\text{HK}(\mathbb{R}^d)$ are truly Riemannian manifolds. 

\subsection{A related approach}
We can pursue a related approach inspired by \citet{gigli2012second}. On a manifold $\man$, one can obtain parallel transport as follows. Let $\gamma(t)$ be a geodesic in $\man$, and let $w_t$ be a vector field along it with $w_0 = w$. It holds that $w_t$ is the parallel transport of $w$ if and only if
\[w_t = d(\exp_{\gamma(0)}(tv))(w)\]
where $v = \dot{\gamma}(0)$ and $d(\cdot)$ is the differential operator. Thus, this gives us a different avenue for pursuing a tractable parallel transport approach along geodesics. If we can compute the differential of the exponential map then we can compute parallel transport.

We'll start with the case of $\mathcal{P}_2(\mathbb{R}^d)$. Given $\mu \in \mathcal{P}_2(\mathbb{R}^d)$ and $u \in L_2(\mu)$, $u:\mathbb{R}^d \rightarrow \mathbb{R}^d$ the exponential map is defined as 
\begin{equation}
    \mathbf{exp}_\mu(u) = (\exp(u))_\#\mu = (\text{Id} + u)_\#\mu
\end{equation}
where $\#$ is the pushforward operator and $\exp$ is the Euclidean exponential map \tls{this was not entirely clear to me - this yields Wasserstein geodesics?} \ndfg{yes it does if we scale it by some $s\in[0,1]$, it traces out a geodesic between the endpoints. It's in Ambrosio gigli Savare in chapter 8 and 12 I think. the $\exp_x(u)$ term is the exponential map on the underlying manifold, which in the flat case is something like $x+u$.}. Note that I will sometimes denote the Wasserstein exponential as $\mathbf{exp}_{\mu, u}$ for convinence. Now pick a regular curve $\mu_t$ with velocity field $v_t$ and an absolutely continuous vector field $u_t$ along it. Further define the curve $t \mapsto \nu_t = \mathbf{exp}_{\mu_t}(u_t)$. \textit{Assuming $\nu_t$ is absolutely continuous}, we'll consider its velocity vector field $w_t$. One should then expect that the differential of the exponential map at the point $\mu_t, u_t$ in the direction $v_t, \dot{u}_t$ is given by $w_t$. We will eventually define the map 
\[j_{\mu, u}: T_u \mathcal{P}_2(\mathbb{R}^d) \times T_u \mathcal{P}_2(\mathbb{R}^d) \rightarrow T_{\exp_{\mu}(u)} \mathcal{P}_2(\mathbb{R}^d) \]
which will represent the \textit{differential} of the exponential map,
\begin{equation}
    j_{\mu, u}(v, w) = d(\mathbf{exp}_{\mu, u}(v, w)).
\end{equation}
Note that the differential of the exponential map takes two arguments, representing instantaneous changes in the measure $\mu$ and the tangent vector $u$. Now, define $\mathbf{T}(t, s, \cdot)$ to be the flow maps of $\mu_t$ and fix some $\varphi \in C^{\infty}(\mathbb{R}^d)$. We'll consider the derivative $\frac{d}{dt}\int\varphi d\nu_t$ shortly, but first observe that
\begin{align}
    \nu_t = \mathbf{exp}_{\mu_t}(u_t) = (\exp(u_t))_\#\mu_t = (\exp(u_t) \circ \mathbf{T}(0,t, \cdot))_\#\mu_0.
\end{align}
Now we can take the derivative,
\begin{align}
    \frac{d}{dt}\int\varphi d\nu_t &= \int \frac{d}{dt}\left[\varphi\circ\exp(u_t) \circ \mathbf{T}(0, t, \cdot)\right] d\mu_0 \\
    &= \int \left\langle (\nabla\varphi) \circ \exp(u_t) \circ \mathbf{T}(0, t, \cdot), \frac{d}{dt} [\exp(u_t) \circ \mathbf{T}(0, t, \cdot)] \right\rangle  d\mu_0.
\end{align} 
Note that in the Euclidean case we have
\begin{align}
    \exp(u_t) \circ \mathbf{T}(0, t, x) = \textbf{T}(0, t, x) + u_t(\textbf{T}(0, t, x)).
\end{align}
Now, Gigli defines the following vector field $J_u[\tilde{u}^1, \tilde{u}^2](t)$ to be the Jacobi field along the geodesic $t \rightarrow \exp_x(tu(x))$ given the initial conditions 
\[J_u[\tilde{u}^1, \tilde{u}^2](0) = \tilde{u}^1(x) \qquad \text{and} \qquad \frac{d}{dt}J_u[\tilde{u}^1, \tilde{u}^2](t)\Big|_{t = 0} = \tilde{u}^2(x).\]
Note that this is a Jacobi field along the \textit{base} manifold, which in our case is $\mathbb{R}^d$. This means we can solve for it quite easily with the Jacobi equation: let $\gamma(t) = x + t\xi$ be a geodesic with velocity $\xi$. Any Jacobi field along $\gamma$ satisfies 
\[\frac{d^2}{dt^2}J(t) + R(J(t), \xi)\xi = 0.\]
Since $\mathbb{R}^d$ is a flat space, its Riemann curvature tensor is identically zero. Thus, the Jacobi equation is satisfied if $J''(t) = 0$. The desired initial conditions imply
\[J_u[\tilde{u}^1, \tilde{u}^2](t) = \tilde{u}^1(x) + t\tilde{u}^2(x).\]
Now lets define the \textit{vector field} $j_{\mu, u}(\tilde{u}^1, \tilde{u}^2)(\exp_x(u(x))) = J_u[\tilde{u}^1, \tilde{u}^2](1) = \tilde{u}^1(x) + \tilde{u}^2(x)$. Proposition 3.13 of \citet{gigli2012second} indicates that
\begin{equation}
    \frac{d}{dt}\left(u_t\left(\mathbf{T}(0, t, x)\right)\right) =(\dot{u}_t)(\mathbf{T}(0, t, x)) \qquad \text{for $\mu_0$ -- a.e. $x$, a.e. $t$}.
\end{equation}
It then follows that
\begin{align}
    \frac{d}{dt}\left(\exp(u_t) \circ \mathbf{T}(0, t, \cdot)\right) &= \frac{d}{dt}(\textbf{T}(0, t, x) + u_t(\textbf{T}(0, t, x)) \\
    &= v_t(\textbf{T}(0, t, x)) + \frac{d}{dt}u_t(\textbf{T}(0, t, x)) \\
    &= v_t(\textbf{T}(0, t, x)) + \dot{u}_t(\textbf{T}(0, t, x)).
\end{align}
It follows that
\[\frac{d}{dt}\left(\exp(u_t) \circ \mathbf{T}(0, t, \cdot)\right) = j_{\mu}(v_t, \dot{u}_t)\Big(\exp(u_t) \circ \mathbf{T}(0, t, \cdot)\Big)\]
which implies
\begin{align}
    \frac{d}{dt}\int \varphi d\nu_t &= \int \bigg\langle \nabla \varphi, j_{\mu, u}(v_t, \dot{u}_t)\bigg\rangle\, d\nu_t \\
    &= \int \left\langle \nabla \varphi, v_t + \dot{u}_t\right\rangle\, d\nu_t \qquad \text{for all } \varphi \in C^{\infty}(\mathbb{R}^d).
\end{align}
This means that the vector field $j_{\mu, u}(v_t, \dot{u}_t)$ satisfies the continuity equation (in the weak sense) for the absolutely continuous curve $\nu_t$. Thus,
\begin{equation}
    w_t = \Pi_{\nu_t}\left(j_{\mu, u}(v_t, \dot{u}_t)\right).
\end{equation}
Specializing this to the Euclidean case yields
\begin{equation}
    w_t = \Pi_{\nu_t}\left(v_t + \dot{u}_t\right).
\end{equation}
Now, we will adapt this strategy for the case of HK. 

\appendix

\section{Technical results}

\begin{restatable}{lemma}{} \label{lemma: orthog jacobi equation}
    For the orthogonal Jacobi equation
    \begin{equation}
        \ddot j_{\perp} +  2\frac{\dot r}{r}\dot j_{\perp} + \frac{2}{r^4}\|q\|_2^2j_{\perp} = 0
    \end{equation}
    we have an equivalent characterization through
    \begin{equation}
        \ddot u_{\perp} + \frac{\|q\|_2^2}{r^4}u_\perp = 0
    \end{equation}
    with $u_\perp = rj_{\perp}$.
\end{restatable}
\noindent \textit{Proof.} We have
\begin{align}
    \dot j_{\perp} &= \frac{r\dot u_{\perp} - u_\perp \dot r}{r^2}
\end{align}
and 
\begin{align}
    \ddot j_{\perp} &= \frac{r^2(\dot r\dot u_{\perp} + r \ddot u_{\perp} - \dot u_\perp \dot r - u_\perp \ddot r) - (r\dot u_{\perp} - u_\perp \dot r)(2r\dot r)}{r^4} \\
    &= \frac{r^2(r \ddot u_{\perp} - u_\perp \ddot r) - (r\dot u_{\perp} - u_\perp \dot r)(2r\dot r)}{r^4} \\
    &= \frac{r^3 \ddot u_{\perp} - r^2u_\perp \ddot r - 2r^2\dot r\dot u_{\perp} + 2ru_\perp \dot r^2}{r^4} \\
    &= \frac{r^2 \ddot u_{\perp} - 2r\dot r\dot u_{\perp}- u_\perp\left(r \ddot r - 2 \dot r^2\right)}{r^3}.
\end{align}
Plugging back into the Jacobi equation yields
\begin{align}
    0 &= \frac{r^2 \ddot u_{\perp} - 2r\dot r\dot u_{\perp}- u_\perp r\ddot r + 2 u_{\perp}\dot r^2}{r^3} + \frac{2\dot rr\dot u_{\perp} - 2u_\perp \dot r^2}{r^3} + \frac{2}{r^5}\|q\|_2^2u_{\perp} \\
    &= \frac{r^2 \ddot u_{\perp} - u_\perp r\ddot r}{r^3} + \frac{2}{r^5}\|q\|_2^2u_{\perp}.
\end{align}
Multiplying through by $r^3$ yields $r^2 \ddot u_{\perp} - u_\perp r\ddot r + \frac{2}{r^2}\|q\|_2^2u_{\perp} = 0$ and dividing by $r^2$ yields $\ddot u_{\perp} - u_\perp \frac{\ddot r}{r} + \frac{2}{r^4}\|q\|_2^2u_{\perp} = 0$. From the geodesic equations we know that $\ddot r = r\|v\|_2^2 = \|q\|^2/r^3$. Thus,
\begin{align}
    \ddot u_{\perp} - u_\perp \frac{\ddot r}{r} + \frac{2}{r^4}\|q\|_2^2u_{\perp} &= \ddot u_{\perp} -  \frac{1}{r^4}\|q\|_2^2u_\perp + \frac{2}{r^4}\|q\|_2^2u_{\perp} \\
    &= \ddot u_{\perp} + \frac{1}{r^4}\|q\|_2^2u_\perp.
\end{align}

\qed{}

\begin{restatable}{lemma}{} \label{lemma: parallel jacobi equation}
    Let $u = b/r$ and suppose that
    \begin{equation}
        \ddot a + 2\frac{\dot r}{r}\dot a + \frac{2}{r^3}\left(\dot b - b\frac{ \dot r}{r}\right) = 0.
    \end{equation}
    and
    \begin{equation}
        \ddot b - \frac{2\dot a}{r}\|q\|_2^2 - \frac{b}{r^4} \|q\|_2^2 = 0.
    \end{equation}
    It follows that
    \begin{equation}
        \dot p + 2\dot u = 0
    \end{equation}
    and 
    \begin{equation}
        (r^2\dot u)' - 2\frac{\|q\|_2^2}{r^2}p = 0
    \end{equation}
    where $u = b/r$ and $p = r^2\dot a$.
\end{restatable}
\noindent \textit{Proof.} The first equation implies
\begin{align}
    0 &= \underbrace{(r^2\ddot a + 2r\dot r\dot a)}_{ = (r^2\dot a)'} + \frac{2}{r}\underbrace{\left(\dot b - b\frac{ \dot r}{r}\right)}_{ = r\dot u} \\
    &= (r^2\dot a)' + 2\dot u.
\end{align}
Taking $p = r^2 \dot a$ finishes the proof of the first equation. The second given equation can be rewritten as 
\begin{align}
    \ddot b  - \frac{b\ddot r}{r} - \frac{2\|q\|_2^2}{r^3}(r^2\dot a) = \ddot b - u\ddot r - \frac{2\|q\|_2^2}{r^3}(r^2\dot a).
\end{align}
since the geodesic equations imply $\ddot r = \|q\|_2^2/r^3.$ Now observe that $\ddot b = \ddot u r + 2\dot u \dot r + u\ddot r$, so 
\begin{align}
    \ddot b - u\ddot r - \frac{2\|q\|_2^2}{r^3}(r^2\dot a) &= \ddot u r  + 2\dot u \dot r - \frac{2\|q\|_2^2}{r^3}(r^2\dot a) \\
    &= \frac{1}{r}(r^2\dot u)' - \frac{2\|q\|_2^2}{r^3}(r^2\dot a).
\end{align}
Thus,
\begin{equation}
    (r^2\dot u)' - \frac{2\|q\|_2^2}{r^2}(r^2\dot a) = 0.
\end{equation}
We obtain the lemma statement by taking $p = r^2\dot a$.  
\qed{}

\begin{restatable}[Time varying oscillator]{lemma}{} \label{lemma: time varying oscillator}
    The ordinary differential equation
    \[\ddot u_{\perp} + \frac{\|q\|_2^2}{r^4}u_\perp = 0\]
    with initial conditions $u_\perp(0)$ and $\dot u_\perp(0)$ satisfies 
    \[u_{\perp}(t) = r(t)\left[\frac{u_{\perp,0}}{r_0}\cos(\|q\|_2\sqrt{2}\tau(t)) + \frac{\dot u_{\perp,0} r_0 - u_{\perp,0} \dot r_0}{\|q\|_2\sqrt{2}}\sin(\|q\|_2\sqrt{2}\tau(t))\right].\]
\end{restatable}
\noindent \textit{Proof.} We want to solve the time-varying harmonic oscillator,
\begin{equation}
    \ddot x + \omega^2(t)x = 0
\end{equation}
where $\omega(t) = \|q\|_2/r^2(t).$ We will start by introducing a complex \textit{ansatz}, and then massaging it to solve our desired ODE. Let
\begin{align}
    \chi(t) = \rho(t)e^{i\theta(t)}
\end{align}
where $\rho >0$ is a real-valued amplitude and $\theta$ is a real-valued phase. We have
\begin{equation}
    \dot \chi = \dot \rho e^{i\theta} + i\dot\theta\rho e^{i\theta}
\end{equation}
and 
\begin{equation}
    \ddot \chi = \left(\ddot \rho  + 2i\dot \theta\dot\rho + i\ddot \theta\rho -\dot\theta^2\rho  \right)e^{i\theta}.
\end{equation}
Plugging into the oscillator and canceling $e^{i\theta}$ yields,
\begin{align}
    (\ddot \rho -\dot\theta^2\rho + \omega^2(t)
    \rho)  + i(2\dot \theta\dot\rho + \ddot \theta\rho) = 0.
\end{align}
This equation is satisfied when both the real and complex parts are null. For the complex component,
\begin{align*}
    2\dot \theta\dot\rho + \ddot \theta\rho &= 0 \\
    \iff (\rho^2\dot\theta)' &= 0 \\
    \iff \rho^2\dot\theta &= C 
\end{align*}
for some $C$ constant in time. So we'll choose $C = \sqrt{2}$, implying $\dot\theta = \sqrt{2}/\rho^2$. Plugging into the real component,
\begin{align*}
    \ddot \rho -\frac{2}{\rho^4}\rho + \omega^2(t)
    \rho &= 0 \\
    \iff\ddot \rho  + \omega^2(t)
    \rho &= \frac{2}{\rho^3}.
\end{align*}
This is an example of an Ermakov-Pinney equation \cite{ermakov1880equations, pinney1950nonlinear}. It turns out that choosing $\rho  = \alpha r$ for some $\alpha > 0$ solves this equation. Verifying,
\begin{align}
    \alpha\ddot r + \frac{\|q\|_2^2}{r^3}\alpha = \frac{2}{\alpha^3r^3}.
\end{align}
The geodesic equations of \cref{eq: geodesic equations on the cone} indicate that $\ddot r = r\|v\|^2_2 = \|q\|_2^2/r^3$ since we have defined $q = r^2v$. Thus,
\begin{align}
    2\frac{\alpha\|q\|_2^2}{r^3} &= \frac{C^2}{\alpha^3r^3} \\
    \implies \alpha &= \frac{1}{\sqrt{\|q\|_2}}
\end{align}
solves the equation. Thus,
\[\rho = \frac{r}{\sqrt{\|q\|_2}}\]
which means
\[\theta(t) = \|q\|_2\sqrt{2}\int_0^t\frac{ds}{r(s)^2} = \|q\|_2\sqrt{2}\tau(t)\]
where $\theta(0) = 0$ by choice. Now the general solution to the original oscillator can be written as 
\begin{align*}
    x(t) &= A\rho(t)\cos\theta(t) + B\rho(t)\sin\theta(t) \\
    &= A\frac{r(t)}{\sqrt{\|q\|_2}}\cos(\|q\|_2\sqrt{2}\tau(t)) + B\frac{r(t)}{\sqrt{\|q\|_2}}\sin(\|q\|_2\sqrt{2}\tau(t)).
\end{align*}
One can solve for $A,B$ with initial conditions $x(0), \dot{x}(0)$ to obtain
\begin{align*}
    A &= \frac{x_0}{r_0}\sqrt{\|q\|_2} \\
    B &= \frac{\dot x_0 r_0 - x_0\dot r_0}{\sqrt{2\|q\|_2}}
\end{align*}
yielding
\begin{equation}
    x(t) = r(t)\left[\frac{x_0}{r_0}\cos(\|q\|_2\sqrt{2}\tau(t)) + \frac{\dot x_0 r_0 - x_0 \dot r_0}{\|q\|_2\sqrt{2}}\sin(\|q\|_2\sqrt{2}\tau(t))\right].
\end{equation}

\qed{}

{\bibliography{main.bib}}

\end{document}
